\subsection{Global approximation with small homotopy tracks}\label{subsec:ar-global-approximation}
\paraid{p-441d3863d070}
In the next theorem,
we combine the local models to construct approximations
through the auxiliary AR
$\AuxiliaryAR=\CompactHyperspace(\BanachAmbientSpace)/\ActingGroup$.
We choose this orbit space for the prescribed error function.
For maps on $\GHSpace$, evaluation at $\MetricSpace{\ComparisonCarrier}{\ComparisonMetric}$
means evaluation at its isometry class $[\ComparisonCarrier,\ComparisonMetric]$.
This convention applies in particular to $\ModelError$ and $\ErrorControl$.

\begin{theorem}\label{thm:global-domination}
\paraid{p-c8631fd0a7ae}
Let
$\ErrorControl \colon \GHSpace\to(0,\infty)$
be continuous.
Then there exist a separable metrizable AR
$\AuxiliaryAR$,
continuous maps
$\ApproximationMap\colon \GHSpace\to \AuxiliaryAR$
and
$\RealizationMap\colon \AuxiliaryAR\to\GHSpace$,
a continuous function
$\ModelError \colon \GHSpace\to[0,\infty)$,
and a continuous homotopy
$\ApproximationHomotopy \colon \GHSpace\times\UnitInterval\to\GHSpace$
with the following properties.
For every
$\MetricSpace{\ComparisonCarrier}{\ComparisonMetric}\in\GHSpace$,
we have
\[
 \ModelError\MetricSpace{\ComparisonCarrier}{\ComparisonMetric}
 <\ErrorControl\MetricSpace{\ComparisonCarrier}{\ComparisonMetric},
\]
and the endpoint identities
\begin{equation}\label{eq:global-domination-1}
 \ApproximationHomotopy (\MetricSpace{\ComparisonCarrier }{\ComparisonMetric },0)=\MetricSpace{\ComparisonCarrier }{\ComparisonMetric },
\end{equation}
and
\begin{equation}\label{eq:global-domination-2}
 \ApproximationHomotopy (\MetricSpace{\ComparisonCarrier }{\ComparisonMetric },1)=\RealizationMap\ApproximationMap\MetricSpace{\ComparisonCarrier }{\ComparisonMetric }.
\end{equation}
For every
$\MetricSpace{\ComparisonCarrier}{\ComparisonMetric}\in\GHSpace$
and every
$\FirstTime,\HomotopyTime\in\UnitInterval$,
we have
\begin{equation}\label{eq:global-track-control}
  \GHDistance(\ApproximationHomotopy (\MetricSpace{\ComparisonCarrier }{\ComparisonMetric },\FirstTime ),\ApproximationHomotopy (\MetricSpace{\ComparisonCarrier }{\ComparisonMetric },\HomotopyTime ))\leq\frac{\abs{\FirstTime -\HomotopyTime }}2 \ModelError \MetricSpace{\ComparisonCarrier }{\ComparisonMetric }.
\end{equation}
\end{theorem}

\begin{proof}
\paraid{p-665396df3bb9}
We construct the local cover and its Banach space,
define the global map,
and then prove continuity and the homotopy estimates.

\ProofStep{Local models and a partition of unity.}\label{step:global-local-models}
\paraid{p-45dee78ed9f8}
For each center
$\MetricSpace{\BaseCarrier}{\MetricSymbol}\in\GHSpace$,
apply \Cref{thm:local-models} with a positive tolerance smaller than
$\ErrorControl\MetricSpace{\BaseCarrier}{\MetricSymbol}$.
By continuity of the model error and of
$\ErrorControl$,
shrink the model domain so that for every space
$\MetricSpace{\ComparisonCarrier}{\ComparisonMetric}$
in that domain the model error is less than
$\ErrorControl\MetricSpace{\ComparisonCarrier}{\ComparisonMetric}$.
Let $\OpenCover$ be the family of these domains.
Then $\OpenCover$ is an open cover of
$\GHSpace$.

\paraid{p-d2972dd59fa1}
Since
$\GHSpace$
is separable and metrizable,
choose a countable locally finite partition of unity
$\{\PartitionWeight _\CoverIndex \}_{\CoverIndex \in\NonnegativeIntegers}$
subordinate to $\OpenCover$
(\Cref{thm:metric-partition-unity}).
For each index
$\CoverIndex$,
fix one of the local models whose domain
$\ModelDomain_\CoverIndex$
contains
$\supp\PartitionWeight_\CoverIndex$.
Different indices may use the same local model.
The supports form a locally finite family.
For each
$\CoverIndex\in\NonnegativeIntegers$,
write the selected local model as
\[
 \mathfrak M_\CoverIndex
 =(\ModelDomain_\CoverIndex,\ModelDimension_\CoverIndex,
   \mathcal L_\CoverIndex,\ModelError_\CoverIndex).
\]
Thus
$\ModelDimension_\CoverIndex\geq1$
is its dimension,
$\mathcal L_\CoverIndex$
is its coordinate-class assignment,
and
$\ModelError_\CoverIndex\colon\ModelDomain_\CoverIndex\to[0,\infty)$
is its error function.
For a compact metric space
$\MetricSpace{\ComparisonCarrier}{\ComparisonMetric}$
whose isometry class belongs to
$\ModelDomain_\CoverIndex$,
choose a pair in
$\mathcal L_\CoverIndex(\ComparisonCarrier,\ComparisonMetric)$.
Write its coordinate map as
$\FeatureCoordinates_{\CoverIndex,\ComparisonCarrier}\colon\ComparisonCarrier\to\RealNumbers^{\ModelDimension_\CoverIndex}$
and its norm as
$\NormSymbol_{\CoverIndex,\ComparisonCarrier}\colon\RealNumbers^{\ModelDimension_\CoverIndex}\to[0,\infty)$.
For
$\ComparisonPoint,\IntegrationPoint\in\ComparisonCarrier$,
write
\[
 \ModelPseudometric_{\CoverIndex,\ComparisonCarrier}(\ComparisonPoint,\IntegrationPoint)
 =\NormSymbol_{\CoverIndex,\ComparisonCarrier}
   (\FeatureCoordinates_{\CoverIndex,\ComparisonCarrier}(\ComparisonPoint)
    -\FeatureCoordinates_{\CoverIndex,\ComparisonCarrier}(\IntegrationPoint)).
\]
The model error satisfies
\[
 \ModelError_\CoverIndex\MetricSpace{\ComparisonCarrier}{\ComparisonMetric}
 =\norm{\ComparisonMetric-\ModelPseudometric_{\CoverIndex,\ComparisonCarrier}}_\infty
 <\ErrorControl\MetricSpace{\ComparisonCarrier}{\ComparisonMetric}.
\]
The zero-weight blocks will be defined separately in Step~\ref{step:global-map}.

\ProofStep{A Banach space with a compact group action.}\label{step:global-orbit-space}
\paraid{p-ab63bc101601}
Form the separable Banach space and compact metrizable group
\begin{equation}\label{eq:banach-sum-group-1}
 \BanachAmbientSpace =\left(\bigoplus_{\CoverIndex \in\NonnegativeIntegers}\BanachAmbientSpace _{\ModelDimension _\CoverIndex }\right)_{\SummableNorm},
\end{equation}
and
\begin{equation}\label{eq:banach-sum-group-2}
 \ActingGroup =\prod_{\CoverIndex \in\NonnegativeIntegers}\OrthogonalGroup(\ModelDimension _\CoverIndex ).
\end{equation}
The norm in
$\BanachAmbientSpace $
is
$\norm{\{\BlockVector _\CoverIndex \}_{\CoverIndex \in\NonnegativeIntegers}}_\BanachAmbientSpace =\sum_{\CoverIndex\in\NonnegativeIntegers} \norm{\BlockVector _\CoverIndex }_\infty$.
Define the induced metric
\[
 \BanachMetric\colon\BanachAmbientSpace\times\BanachAmbientSpace\to[0,\infty)
\]
as follows.
For every
$\FirstVector,\SecondVector\in\BanachAmbientSpace$,
set
\[
 \BanachMetric(\FirstVector,\SecondVector)=\norm{\FirstVector-\SecondVector}_\BanachAmbientSpace.
\]
The group
$\ActingGroup$
acts on each summand by linear isometries.
We prove that the action map
\[
 \ActingGroup\times\BanachAmbientSpace\to\BanachAmbientSpace,
 \qquad
 (\GroupElement,\DirectSumVector)\longmapsto\GroupElement\DirectSumVector
\]
is continuous.
Indeed,
if
$\GroupElement _\SequenceIndex \to \GroupElement $
in the product group and
$\DirectSumVector =\{\BlockVector _\CoverIndex \}\in \BanachAmbientSpace $,
then for each
$\TruncationSize \geq1$,
\[
 \norm{\GroupElement _\SequenceIndex \DirectSumVector -\GroupElement \DirectSumVector }_\BanachAmbientSpace
 \leq\sum_{\CoverIndex =0}^\TruncationSize \norm{(\GroupElement _\SequenceIndex )_\CoverIndex \BlockVector _\CoverIndex -\GroupElement _\CoverIndex \BlockVector _\CoverIndex }_\infty
        +2\sum_{\CoverIndex >\TruncationSize }\norm{\BlockVector _\CoverIndex }_\infty.
\]
First choose
$\TruncationSize $
to control the tail,
and then use continuity on the finitely many remaining blocks.
For varying
$\ApproximatingDirectSumVector _\SequenceIndex \to \DirectSumVector $,
the inequality
\[
 \norm{\GroupElement _\SequenceIndex \ApproximatingDirectSumVector _\SequenceIndex -\GroupElement \DirectSumVector }_\BanachAmbientSpace \leq\norm{\ApproximatingDirectSumVector _\SequenceIndex -\DirectSumVector }_\BanachAmbientSpace +\norm{\GroupElement _\SequenceIndex \DirectSumVector -\GroupElement \DirectSumVector }_\BanachAmbientSpace\longrightarrow0
\]
proves continuity of the action map on the product
$\ActingGroup\times\BanachAmbientSpace$.

\ProofStep{The global map and its error.}\label{step:global-map}
\paraid{p-4b598342fd3f}
Fix a compact metric space
$\MetricSpace{\ComparisonCarrier}{\ComparisonMetric}$
representing an input class in
$\GHSpace$.
Define the finite set of active indices for this input by
\[
 \mathcal I(\ComparisonCarrier,\ComparisonMetric)
 =\{\CoverIndex\in\NonnegativeIntegers\mid
 \PartitionWeight_\CoverIndex\MetricSpace{\ComparisonCarrier}{\ComparisonMetric}>0\}.
\]
Finiteness follows from local finiteness of the partition of unity.
For each index with
$\PartitionWeight_\CoverIndex\MetricSpace{\ComparisonCarrier}{\ComparisonMetric}>0$,
choose a representative coordinate pair
$(\FeatureCoordinates_{\CoverIndex,\ComparisonCarrier},\NormSymbol_{\CoverIndex,\ComparisonCarrier})$
from the coordinate class of the assigned local model.
For every $\CoverIndex\in\NonnegativeIntegers$ and
$\ComparisonPoint\in\ComparisonCarrier$, define
\[
 B_{\CoverIndex,\ComparisonCarrier}(\ComparisonPoint)=
 \begin{cases}
 \PartitionWeight_\CoverIndex\MetricSpace{\ComparisonCarrier}{\ComparisonMetric}
 \NormEmbedding_{\NormSymbol_{\CoverIndex,\ComparisonCarrier}}
 (\FeatureCoordinates_{\CoverIndex,\ComparisonCarrier}(\ComparisonPoint))
 &\text{if }\PartitionWeight_\CoverIndex\MetricSpace{\ComparisonCarrier}{\ComparisonMetric}>0,\\
 0&\text{otherwise}.
 \end{cases}
\]
The zero in the second case is the zero vector of
$\BanachAmbientSpace_{\ModelDimension_\CoverIndex}$.
Using these blocks,
define the map
$\JointFeatureMap_\ComparisonCarrier\colon\ComparisonCarrier\to\BanachAmbientSpace$
as follows.
For every
$\ComparisonPoint\in\ComparisonCarrier$,
put
\begin{equation}\label{eq:joint-feature-map-1}
 \JointFeatureMap_\ComparisonCarrier(\ComparisonPoint)=\{B_{\CoverIndex,\ComparisonCarrier}(\ComparisonPoint)\}_{\CoverIndex\in\NonnegativeIntegers},
\end{equation}
and denote its image by
\begin{equation}\label{eq:joint-feature-map-2}
 \CompactFeatureImage _\ComparisonCarrier =\JointFeatureMap _\ComparisonCarrier (\ComparisonCarrier ).
\end{equation}
For every point of
$\ComparisonCarrier$,
all nonzero blocks have indices in the same finite set
$\mathcal I(\ComparisonCarrier,\ComparisonMetric)$.
Thus
$\JointFeatureMap _\ComparisonCarrier $
is continuous and
$\CompactFeatureImage _\ComparisonCarrier $
is nonempty and compact.
Equation \eqref{eq:dual-isometry} and \ref{item:model-bound} of \Cref{thm:local-models}
imply that for every
$\ComparisonPoint\in\ComparisonCarrier$
\begin{equation}\label{eq:global-feature-bound}
  \norm{\JointFeatureMap _\ComparisonCarrier (\ComparisonPoint )}_\BanachAmbientSpace \leq2\diam\MetricSpace{\ComparisonCarrier }{\ComparisonMetric }.
\end{equation}

\paraid{p-b86cdaebafd5}
If the local coordinate pairs are changed as in \eqref{eq:frame-action} by
$\GroupElement=\{\OrthogonalChange_\CoverIndex\}_\CoverIndex\in\ActingGroup$,
then \eqref{eq:dual-equivariance} transforms the joint point map into
$\GroupElement\circ\JointFeatureMap_\ComparisonCarrier$
and its image into
$\GroupElement\CompactFeatureImage_\ComparisonCarrier$.
If
$\InputIsometry\colon\MetricSpace{\ComparisonCarrier}{\ComparisonMetric}
 \to\MetricSpace{\ComparisonCarrier_0}{\ComparisonMetric_0}$
is an isometry,
then
\eqref{eq:input-isometry-1},
\eqref{eq:input-isometry-2},
and \eqref{eq:dual-equivariance}
imply that for some
$\GroupElement\in\ActingGroup$,
\[
 \JointFeatureMap_{\ComparisonCarrier_0}\circ\InputIsometry
 =\GroupElement\circ\JointFeatureMap_\ComparisonCarrier.
\]
Thus both changes leave the orbit of the compact image unchanged.
Consequently,
we can define
\[
 \ApproximationMap\colon \GHSpace\to \AuxiliaryAR=\CompactHyperspace(\BanachAmbientSpace )/\ActingGroup
\]
by
\[
 \ApproximationMap\MetricSpace{\ComparisonCarrier }{\ComparisonMetric }=[\CompactFeatureImage _\ComparisonCarrier ].
\]
The value is independent of the representative of the input isometry class.
By \Cref{lem:hyperspace-orbit-ar},
$\AuxiliaryAR$
is a separable metrizable AR.
By \Cref{lem:hyperspace-orbit-realization},
we obtain a
$1$-Lipschitz map
$\RealizationMap\colon \AuxiliaryAR\to\GHSpace$.

\paraid{p-1264f3c8b4b3}
Figure~\ref{fig:global-models} shows how the local maps are combined.
\begin{figure}[H]
\centering
\begin{tikzpicture}[>=Stealth,font=\small]
\node (input) at (0,0) {$\ComparisonPoint\in\ComparisonCarrier$};
\node (first) at (4,1) {$\PartitionWeight_{\CoverIndex_0}\MetricSpace{\ComparisonCarrier}{\ComparisonMetric}
 \NormEmbedding_{\NormSymbol_{\CoverIndex_0,\ComparisonCarrier}}
 (\FeatureCoordinates_{\CoverIndex_0,\ComparisonCarrier}(\ComparisonPoint))$};
\node (second) at (4,-1) {$\PartitionWeight_{\CoverIndex_1}\MetricSpace{\ComparisonCarrier}{\ComparisonMetric}
 \NormEmbedding_{\NormSymbol_{\CoverIndex_1,\ComparisonCarrier}}
 (\FeatureCoordinates_{\CoverIndex_1,\ComparisonCarrier}(\ComparisonPoint))$};
\node at (4,0) {$\vdots$};
\node (joint) at (8,0) {$\JointFeatureMap_\ComparisonCarrier(\ComparisonPoint)\in\BanachAmbientSpace$};
\draw[->] (input)--(first);
\draw[->] (input)--(second);
\draw[->] (first)--(joint);
\draw[->] (second)--(joint);
\end{tikzpicture}
\par\medskip
\begin{tikzcd}[column sep=large]
 \GHSpace \arrow[r,"\ApproximationMap"] &
 \AuxiliaryAR=\CompactHyperspace(\BanachAmbientSpace)/\ActingGroup
 \arrow[r,"\RealizationMap"] & \GHSpace
\end{tikzcd}
\[
 \RealizationMap\ApproximationMap\simeq\IdentityMap_\GHSpace
 \quad\text{via }\ApproximationHomotopy.
\]
\caption{For a fixed compact metric space
$\MetricSpace{\ComparisonCarrier}{\ComparisonMetric}$,
the active blocks use the same point
$\ComparisonPoint$.
Their joint image is
$\CompactFeatureImage_\ComparisonCarrier=\JointFeatureMap_\ComparisonCarrier(\ComparisonCarrier)$.
Changing the chosen coordinate pairs as in \eqref{eq:frame-action} acts on this image by
$\ActingGroup$,
so
$\ApproximationMap\MetricSpace{\ComparisonCarrier}{\ComparisonMetric}
 =[\CompactFeatureImage_\ComparisonCarrier]$.
The displayed composite is joined to the identity by the homotopy in Step~\ref{step:global-homotopy}
with the track estimate \eqref{eq:global-track-control}.}
\label{fig:global-models}
\end{figure}


\paraid{p-b45630ebc16d}
For
$\ComparisonPoint,\IntegrationPoint\in\ComparisonCarrier$,
define the pseudometric induced by
$\JointFeatureMap_\ComparisonCarrier$
by
\begin{equation}\label{eq:global-pseudometric}
  \ModelPseudometric _\ComparisonCarrier (\ComparisonPoint ,\IntegrationPoint )=\norm{\JointFeatureMap _\ComparisonCarrier (\ComparisonPoint )-\JointFeatureMap _\ComparisonCarrier (\IntegrationPoint )}_\BanachAmbientSpace.
\end{equation}
By \eqref{eq:dual-isometry},
we have
\begin{equation}\label{eq:global-pseudometric-sum}
 \ModelPseudometric_\ComparisonCarrier(\ComparisonPoint,\IntegrationPoint)
 =\sum_{\substack{\CoverIndex\in\NonnegativeIntegers\\\PartitionWeight_\CoverIndex\MetricSpace{\ComparisonCarrier}{\ComparisonMetric}>0}} \PartitionWeight _\CoverIndex \MetricSpace{\ComparisonCarrier }{\ComparisonMetric }\ModelPseudometric _{\CoverIndex ,\ComparisonCarrier }(\ComparisonPoint ,\IntegrationPoint ).
\end{equation}
Every block in \eqref{eq:joint-feature-map-1} is evaluated at the same point
$\ComparisonPoint\in\ComparisonCarrier$.
Define the error by
\begin{equation}\label{eq:global-error}
  \ModelError \MetricSpace{\ComparisonCarrier }{\ComparisonMetric }=\norm{\ComparisonMetric -\ModelPseudometric _\ComparisonCarrier }_\infty.
\end{equation}
Since the weights sum to one,
\eqref{eq:global-pseudometric-sum} implies
\begin{equation}\label{eq:global-error-bound}
 \ModelError\MetricSpace{\ComparisonCarrier}{\ComparisonMetric}
 \leq
 \sum_{\CoverIndex\in \mathcal I(\ComparisonCarrier,\ComparisonMetric)}\PartitionWeight _\CoverIndex \MetricSpace{\ComparisonCarrier }{\ComparisonMetric }\ModelError _\CoverIndex \MetricSpace{\ComparisonCarrier }{\ComparisonMetric }
 <
 \sum_{\CoverIndex\in \mathcal I(\ComparisonCarrier,\ComparisonMetric)}\PartitionWeight _\CoverIndex \MetricSpace{\ComparisonCarrier }{\ComparisonMetric }
 \ErrorControl \MetricSpace{\ComparisonCarrier }{\ComparisonMetric }
=\ErrorControl \MetricSpace{\ComparisonCarrier }{\ComparisonMetric }.
\end{equation}





\ProofStep{Continuity along common correspondences.}\label{step:global-continuity}
\paraid{p-613bbdb06ba1}
We next prove continuity of
$\ApproximationMap$
and
$\ModelError$
by estimating the maps
$\JointFeatureMap_\ComparisonCarrier$
along common correspondences.
Local finiteness reduces the argument to finitely many blocks.
We treat separately the blocks whose weights have positive limits
and those whose weights tend to
$0$.
We will also use the estimate \eqref{eq:global-alignment} to prove continuity of
$\ApproximationHomotopy$.

\ProofSubstep{A common correspondence and finitely many blocks.}\label{step:global-common-correspondence}
\paraid{p-527721648b5d}
Let
$\MetricSpace{\ComparisonCarrier _\SequenceIndex }{\ComparisonMetric _\SequenceIndex}\to\MetricSpace{\ComparisonCarrier }{\ComparisonMetric }$
in
$\GHSpace$.
By \Cref{lem:common-gh-embedding},
we may regard all the carriers as compact subsets of one compact metric space
$\MetricSpace{\AmbientCarrier}{\AmbientMetric}$
such that
$\HausdorffDistance{\AmbientMetric}
 (\ComparisonCarrier_\SequenceIndex,\ComparisonCarrier)\to0$.
Choose correspondences
$\Correspondence _\SequenceIndex \subset \ComparisonCarrier _\SequenceIndex \times \ComparisonCarrier $
such that
\[
 \AmbientDisplacement_\SequenceIndex
 =\sup_{(\ComparisonPoint_\SequenceIndex,\ComparisonPoint)
          \in\Correspondence_\SequenceIndex}
   \AmbientMetric(\ComparisonPoint_\SequenceIndex,\ComparisonPoint)
 \longrightarrow0.
\]
We use these same correspondences for every block.
Local finiteness of the supports provides an open neighborhood
$\OpenNeighborhood$
of
$\MetricSpace{\ComparisonCarrier}{\ComparisonMetric}$
such that
\[
 \ActiveModelIndices
 =\{\CoverIndex\in\NonnegativeIntegers\mid
       \OpenNeighborhood\cap\supp\PartitionWeight_\CoverIndex
       \ne\emptyset\}
 \quad\text{is finite}.
\]
For all sufficiently large
$\SequenceIndex$,
we have
$\MetricSpace{\ComparisonCarrier_\SequenceIndex}{\ComparisonMetric_\SequenceIndex}
 \in\OpenNeighborhood$.
Thus,
for every
$\CoverIndex\in\NonnegativeIntegers\setminus\ActiveModelIndices$,
\[
 \PartitionWeight_\CoverIndex\MetricSpace{\ComparisonCarrier}{\ComparisonMetric}
 =\PartitionWeight_\CoverIndex
   \MetricSpace{\ComparisonCarrier_\SequenceIndex}{\ComparisonMetric_\SequenceIndex}
 =0.
\]
Partition the finite set
$\ActiveModelIndices$
according to the limiting weights by putting
\[
 I_+=\{\CoverIndex\in\ActiveModelIndices\mid
 \PartitionWeight_\CoverIndex\MetricSpace{\ComparisonCarrier}{\ComparisonMetric}>0\}
\]
and
$I_-=\ActiveModelIndices\setminus I_+$.
Namely,
\[
 I_-=\{\CoverIndex\in\ActiveModelIndices\mid
 \PartitionWeight_\CoverIndex\MetricSpace{\ComparisonCarrier}{\ComparisonMetric}=0\}.
\]

\ProofSubstep{Blocks with positive limiting weights.}\label{step:global-positive-blocks}
\paraid{p-5b0d3e37abfc}
For each
$\CoverIndex\in I_+$,
continuity of
$\PartitionWeight_\CoverIndex$
implies
\[
 \PartitionWeight_\CoverIndex
 \MetricSpace{\ComparisonCarrier_\SequenceIndex}{\ComparisonMetric_\SequenceIndex}
 \longrightarrow
 \PartitionWeight_\CoverIndex\MetricSpace{\ComparisonCarrier}{\ComparisonMetric}>0
 \quad\text{as }\SequenceIndex\to\infty.
\]
Since
$I_+$
is finite,
we ignore finitely many initial terms so that
$\PartitionWeight_\CoverIndex
 \MetricSpace{\ComparisonCarrier_\SequenceIndex}{\ComparisonMetric_\SequenceIndex}>0$
for every
$\CoverIndex\in I_+$
and every remaining index
$\SequenceIndex$.
Thus,
for each
$\CoverIndex\in I_+$,
the corresponding local model is defined at both
$\MetricSpace{\ComparisonCarrier_\SequenceIndex}{\ComparisonMetric_\SequenceIndex}$
and
$\MetricSpace{\ComparisonCarrier}{\ComparisonMetric}$.
We now examine the continuity of the factor
$\NormEmbedding_{\NormSymbol_{\CoverIndex,\ComparisonCarrier}}
 \circ\FeatureCoordinates_{\CoverIndex,\ComparisonCarrier}$
in the definition of
$B_{\CoverIndex,\ComparisonCarrier}$
as both the norm and the coordinate map vary.
By \ref{item:model-continuity} of \Cref{thm:local-models},
choose coordinate pairs for each
$\CoverIndex\in I_+$
such that
\begin{equation}\label{eq:global-block-norm-convergence}
 \sup_{\SpherePoint\in\UnitSphere^{\ModelDimension_\CoverIndex-1}}
 \bigl|\NormSymbol_{\CoverIndex,\ComparisonCarrier_\SequenceIndex}(\SpherePoint)
       -\NormSymbol_{\CoverIndex,\ComparisonCarrier}(\SpherePoint)\bigr|
 \longrightarrow0
\end{equation}
and
\begin{equation}\label{eq:global-block-coordinate-convergence}
 \sup_{(\ComparisonPoint_\SequenceIndex,\ComparisonPoint)
          \in\Correspondence_\SequenceIndex}
 \norm{\FeatureCoordinates_{\CoverIndex,\ComparisonCarrier_\SequenceIndex}
          (\ComparisonPoint_\SequenceIndex)
       -\FeatureCoordinates_{\CoverIndex,\ComparisonCarrier}
          (\ComparisonPoint)}_2
\longrightarrow0.
\end{equation}
Fix
$\CoverIndex\in I_+$.
For every sufficiently large
$\SequenceIndex$,
the supremum in \eqref{eq:global-block-coordinate-convergence}
is at most
$1$.
Fix such an index
$\SequenceIndex$
and take an arbitrary point
$\ComparisonPoint_\SequenceIndex\in\ComparisonCarrier_\SequenceIndex$.
Since the first coordinate projection
$\Correspondence_\SequenceIndex\to\ComparisonCarrier_\SequenceIndex$
is surjective,
choose
$\ComparisonPoint\in\ComparisonCarrier$
such that
$(\ComparisonPoint_\SequenceIndex,\ComparisonPoint)\in\Correspondence_\SequenceIndex$.
Applying the triangle inequality to this pair,
we obtain
\begin{equation}\label{eq:global-block-pointwise-bound}
 \begin{aligned}
 \norm{\FeatureCoordinates_{\CoverIndex,\ComparisonCarrier_\SequenceIndex}
       (\ComparisonPoint_\SequenceIndex)}_2
 &\leq
 \norm{\FeatureCoordinates_{\CoverIndex,\ComparisonCarrier_\SequenceIndex}
       (\ComparisonPoint_\SequenceIndex)
       -\FeatureCoordinates_{\CoverIndex,\ComparisonCarrier}(\ComparisonPoint)}_2
 +\norm{\FeatureCoordinates_{\CoverIndex,\ComparisonCarrier}(\ComparisonPoint)}_2\\
 &\leq1+\max_{\IntegrationPoint\in\ComparisonCarrier}
       \norm{\FeatureCoordinates_{\CoverIndex,\ComparisonCarrier}(\IntegrationPoint)}_2.
 \end{aligned}
\end{equation}
Taking the supremum over all
$\ComparisonPoint_\SequenceIndex\in\ComparisonCarrier_\SequenceIndex$
in \eqref{eq:global-block-pointwise-bound},
we obtain
\begin{equation}\label{eq:global-block-coordinate-bound}
 \sup_{\ComparisonPoint_\SequenceIndex\in\ComparisonCarrier_\SequenceIndex}
 \norm{\FeatureCoordinates_{\CoverIndex,\ComparisonCarrier_\SequenceIndex}
       (\ComparisonPoint_\SequenceIndex)}_2
 \leq1+\max_{\IntegrationPoint\in\ComparisonCarrier}
       \norm{\FeatureCoordinates_{\CoverIndex,\ComparisonCarrier}(\IntegrationPoint)}_2.
\end{equation}
Put
\[
 \Radius_\CoverIndex
 =1+\max_{\IntegrationPoint\in\ComparisonCarrier}
       \norm{\FeatureCoordinates_{\CoverIndex,\ComparisonCarrier}(\IntegrationPoint)}_2
 <\infty.
\]
The finiteness follows from continuity and compactness.
By \eqref{eq:global-block-coordinate-bound},
for all sufficiently large
$\SequenceIndex$,
the images of
$\FeatureCoordinates_{\CoverIndex,\ComparisonCarrier_\SequenceIndex}$
and
$\FeatureCoordinates_{\CoverIndex,\ComparisonCarrier}$
lie in the closed Euclidean ball
\[
 \overline{\MetricBall}(0,\Radius_\CoverIndex;\norm{\cdot}_2)
 =\{\FirstVector\in\RealNumbers^{\ModelDimension_\CoverIndex}\mid
     \norm{\FirstVector}_2\leq\Radius_\CoverIndex\}.
\]
For
$(\ComparisonPoint_\SequenceIndex,\ComparisonPoint)
 \in\Correspondence_\SequenceIndex$,
we separate the change of norm from the change of coordinates by the estimate
\begin{equation}\label{eq:global-block-embedding-comparison}
 \begin{aligned}
 &\Bigl\|\NormEmbedding_{\NormSymbol_{\CoverIndex,\ComparisonCarrier_\SequenceIndex}}
   \bigl(\FeatureCoordinates_{\CoverIndex,\ComparisonCarrier_\SequenceIndex}
      (\ComparisonPoint_\SequenceIndex)\bigr)
 -\NormEmbedding_{\NormSymbol_{\CoverIndex,\ComparisonCarrier}}
   \bigl(\FeatureCoordinates_{\CoverIndex,\ComparisonCarrier}
      (\ComparisonPoint)\bigr)\Bigr\|_\infty\\
 &\quad\leq
 \Bigl\|\bigl(\NormEmbedding_{\NormSymbol_{\CoverIndex,\ComparisonCarrier_\SequenceIndex}}
             -\NormEmbedding_{\NormSymbol_{\CoverIndex,\ComparisonCarrier}}\bigr)
   \bigl(\FeatureCoordinates_{\CoverIndex,\ComparisonCarrier_\SequenceIndex}
      (\ComparisonPoint_\SequenceIndex)\bigr)\Bigr\|_\infty
 +
 \NormSymbol_{\CoverIndex,\ComparisonCarrier}
   \bigl(\FeatureCoordinates_{\CoverIndex,\ComparisonCarrier_\SequenceIndex}
      (\ComparisonPoint_\SequenceIndex)
       -\FeatureCoordinates_{\CoverIndex,\ComparisonCarrier}(\ComparisonPoint)\bigr)\\
 &\quad\leq
 \sup_{\FirstVector\in\overline{\MetricBall}(0,\Radius_\CoverIndex;\norm{\cdot}_2)}
 \Bigl\|\bigl(\NormEmbedding_{\NormSymbol_{\CoverIndex,\ComparisonCarrier_\SequenceIndex}}
             -\NormEmbedding_{\NormSymbol_{\CoverIndex,\ComparisonCarrier}}\bigr)
   (\FirstVector)\Bigr\|_\infty
 +
 \NormSymbol_{\CoverIndex,\ComparisonCarrier}
   \bigl(\FeatureCoordinates_{\CoverIndex,\ComparisonCarrier_\SequenceIndex}
      (\ComparisonPoint_\SequenceIndex)
       -\FeatureCoordinates_{\CoverIndex,\ComparisonCarrier}(\ComparisonPoint)\bigr).
 \end{aligned}
\end{equation}
The first inequality in \eqref{eq:global-block-embedding-comparison}
follows from the triangle inequality and \eqref{eq:dual-isometry}.
The second follows from \eqref{eq:global-block-coordinate-bound}.
By \eqref{eq:global-block-norm-convergence}
and \eqref{eq:dual-embedding-uniform-on-bounded},
we obtain
\begin{equation}\label{eq:global-block-embedding-uniform-convergence}
 \sup_{\FirstVector\in\overline{\MetricBall}(0,\Radius_\CoverIndex;\norm{\cdot}_2)}
 \Bigl\|\bigl(\NormEmbedding_{\NormSymbol_{\CoverIndex,\ComparisonCarrier_\SequenceIndex}}
             -\NormEmbedding_{\NormSymbol_{\CoverIndex,\ComparisonCarrier}}\bigr)
   (\FirstVector)\Bigr\|_\infty
 \longrightarrow0.
\end{equation}
Since the fixed norm
$\NormSymbol_{\CoverIndex,\ComparisonCarrier}$
is continuous on the compact unit sphere,
the constant
\[
 \NormComparisonBound_\CoverIndex
 =\max_{\SpherePoint\in\UnitSphere^{\ModelDimension_\CoverIndex-1}}
       \NormSymbol_{\CoverIndex,\ComparisonCarrier}(\SpherePoint)
\]
is finite.
For every
$\FirstVector\in\RealNumbers^{\ModelDimension_\CoverIndex}$,
homogeneity implies
\[
 \NormSymbol_{\CoverIndex,\ComparisonCarrier}(\FirstVector)
 \leq\NormComparisonBound_\CoverIndex\norm{\FirstVector}_2.
\]
Thus,
\eqref{eq:global-block-coordinate-convergence} implies
\begin{equation}\label{eq:global-block-coordinate-norm-convergence}
 \begin{aligned}
 &\sup_{(\ComparisonPoint_\SequenceIndex,\ComparisonPoint)
          \in\Correspondence_\SequenceIndex}
 \NormSymbol_{\CoverIndex,\ComparisonCarrier}
   \bigl(\FeatureCoordinates_{\CoverIndex,\ComparisonCarrier_\SequenceIndex}
      (\ComparisonPoint_\SequenceIndex)
       -\FeatureCoordinates_{\CoverIndex,\ComparisonCarrier}(\ComparisonPoint)\bigr)\\
 &\quad\leq\NormComparisonBound_\CoverIndex
 \sup_{(\ComparisonPoint_\SequenceIndex,\ComparisonPoint)
          \in\Correspondence_\SequenceIndex}
 \norm{\FeatureCoordinates_{\CoverIndex,\ComparisonCarrier_\SequenceIndex}
          (\ComparisonPoint_\SequenceIndex)
       -\FeatureCoordinates_{\CoverIndex,\ComparisonCarrier}
          (\ComparisonPoint)}_2
 \longrightarrow0.
 \end{aligned}
\end{equation}
By \eqref{eq:global-block-embedding-comparison},
\eqref{eq:global-block-embedding-uniform-convergence},
and \eqref{eq:global-block-coordinate-norm-convergence},
for every
$\CoverIndex\in I_+$,
the embedded coordinate maps are uniformly bounded and converge
uniformly along
$\Correspondence_\SequenceIndex$.
By continuity of the weights
$\PartitionWeight_\CoverIndex$
and the definition of the blocks
$B_{\CoverIndex,\ComparisonCarrier_\SequenceIndex}$
and
$B_{\CoverIndex,\ComparisonCarrier}$
in Step~\ref{step:global-map},
we obtain
\begin{equation}\label{eq:global-positive-weight-block}
 \sup_{(\ComparisonPoint_\SequenceIndex,\ComparisonPoint)
          \in\Correspondence_\SequenceIndex}
 \norm{B_{\CoverIndex,\ComparisonCarrier_\SequenceIndex}
           (\ComparisonPoint_\SequenceIndex)
       -B_{\CoverIndex,\ComparisonCarrier}(\ComparisonPoint)}_\infty
 \longrightarrow0.
\end{equation}

\ProofSubstep{Blocks with vanishing weights.}\label{step:global-vanishing-blocks}
\paraid{p-50b148b589f3}
Fix
$\CoverIndex\in I_-$.
By continuity of
$\PartitionWeight_\CoverIndex$,
we have
$\PartitionWeight_\CoverIndex
 \MetricSpace{\ComparisonCarrier_\SequenceIndex}{\ComparisonMetric_\SequenceIndex}
 \to\PartitionWeight_\CoverIndex\MetricSpace{\ComparisonCarrier}{\ComparisonMetric}=0$.
The definition in Step~\ref{step:global-map} implies
$B_{\CoverIndex,\ComparisonCarrier}=0$.
If the weight at
$\MetricSpace{\ComparisonCarrier_\SequenceIndex}{\ComparisonMetric_\SequenceIndex}$
is zero,
then the corresponding block is zero by its definition in Step~\ref{step:global-map}.
If the weight is positive,
then the local model is defined,
and \ref{item:model-bound} of \Cref{thm:local-models}
and \eqref{eq:dual-isometry} apply.
In both cases,
\begin{equation}\label{eq:vanishing-weight-block}
 \sup_{\ComparisonPoint_\SequenceIndex\in\ComparisonCarrier_\SequenceIndex}
 \norm{B_{\CoverIndex,\ComparisonCarrier_\SequenceIndex}(\ComparisonPoint_\SequenceIndex)}_\infty
 \leq2\PartitionWeight _\CoverIndex \MetricSpace{\ComparisonCarrier _\SequenceIndex }{\ComparisonMetric _\SequenceIndex}\diam\MetricSpace{\ComparisonCarrier _\SequenceIndex }{\ComparisonMetric _\SequenceIndex}\to0.
\end{equation}
Indeed,
the diameters are bounded along the convergent sequence.
Thus these blocks converge uniformly to
$0$
for any choices of coordinate pairs.

\ProofSubstep{Continuity of the global map and its error.}\label{step:global-map-continuity}
\paraid{p-f79a92fe8ba3}
Use the coordinate pairs chosen in Step~\ref{step:global-positive-blocks} on
$I_+$
and arbitrary coordinate pairs for the blocks with positive weights in
$I_-$.
Put
\[
 \FeatureDisplacement _\SequenceIndex =\sup_{(\ComparisonPoint _\SequenceIndex ,\ComparisonPoint )\in \Correspondence _\SequenceIndex }\norm{\JointFeatureMap _{\ComparisonCarrier _\SequenceIndex }(\ComparisonPoint _\SequenceIndex )-\JointFeatureMap _\ComparisonCarrier (\ComparisonPoint )}_\BanachAmbientSpace ,
\]
so that
$\FeatureDisplacement_\SequenceIndex$
measures the displacement of the joint feature maps along
$\Correspondence_\SequenceIndex$.
All blocks outside
$\ActiveModelIndices$
vanish for sufficiently large
$\SequenceIndex$.
Since the norm of
$\BanachAmbientSpace$
is the sum of the block norms and
$\ActiveModelIndices$
is finite,
\eqref{eq:global-positive-weight-block} and \eqref{eq:vanishing-weight-block}
imply
\begin{equation}\label{eq:global-alignment}
 \FeatureDisplacement_\SequenceIndex
 \leq\sum_{\CoverIndex\in\ActiveModelIndices}
 \sup_{(\ComparisonPoint_\SequenceIndex,\ComparisonPoint)
          \in\Correspondence_\SequenceIndex}
 \norm{B_{\CoverIndex,\ComparisonCarrier_\SequenceIndex}
           (\ComparisonPoint_\SequenceIndex)
       -B_{\CoverIndex,\ComparisonCarrier}(\ComparisonPoint)}_\infty
\longrightarrow0.
\end{equation}
Since both coordinate projections of
$\Correspondence_\SequenceIndex$
are onto,
the definition of
$\FeatureDisplacement_\SequenceIndex$
implies
\[
 \HausdorffDistance{\BanachMetric}
 (\CompactFeatureImage_{\ComparisonCarrier_\SequenceIndex},
  \CompactFeatureImage_\ComparisonCarrier)
 \leq\FeatureDisplacement_\SequenceIndex\longrightarrow0.
\]
Here the representatives
$\CompactFeatureImage_{\ComparisonCarrier_\SequenceIndex}$
and
$\CompactFeatureImage_\ComparisonCarrier$
are constructed using,
for each
$\CoverIndex\in I_+$,
the coordinate pairs guaranteed by
\ref{item:model-continuity} of \Cref{thm:local-models}
to satisfy \eqref{eq:global-block-norm-convergence}
and \eqref{eq:global-block-coordinate-convergence}.
Taking the identity element of
$\ActingGroup$
in \eqref{eq:orbit-metric},
we obtain
\[
 \OrbitMetric([\CompactFeatureImage_{\ComparisonCarrier_\SequenceIndex}],
             [\CompactFeatureImage_\ComparisonCarrier])
 \leq\HausdorffDistance{\BanachMetric}
 (\CompactFeatureImage_{\ComparisonCarrier_\SequenceIndex},
  \CompactFeatureImage_\ComparisonCarrier)
 \longrightarrow0.
\]
This proves continuity of
$\ApproximationMap$.

\paraid{p-39150d57f5ed}
We finally prove continuity of
$\ModelError$.
For
$(\ComparisonPoint _\SequenceIndex ,\ComparisonPoint ),(\IntegrationPoint _\SequenceIndex ,\IntegrationPoint )\in \Correspondence _\SequenceIndex $,
the triangle inequalities in
$\BanachAmbientSpace$
and
$\MetricSpace{\AmbientCarrier}{\AmbientMetric}$
imply
\[
 |\ModelPseudometric _{\ComparisonCarrier _\SequenceIndex }(\ComparisonPoint _\SequenceIndex ,\IntegrationPoint _\SequenceIndex )-\ModelPseudometric _\ComparisonCarrier (\ComparisonPoint ,\IntegrationPoint )|\leq2\FeatureDisplacement _\SequenceIndex ,
\]
and
\[
 |\ComparisonMetric _\SequenceIndex(\ComparisonPoint _\SequenceIndex ,\IntegrationPoint _\SequenceIndex )-\ComparisonMetric (\ComparisonPoint ,\IntegrationPoint )|\leq2\AmbientDisplacement _\SequenceIndex .
\]
By \eqref{eq:pseudometric-uniform-comparison},
we obtain
\[
 |\ModelError \MetricSpace{\ComparisonCarrier _\SequenceIndex }{\ComparisonMetric _\SequenceIndex}-\ModelError \MetricSpace{\ComparisonCarrier }{\ComparisonMetric }|
 \leq2\FeatureDisplacement _\SequenceIndex +2\AmbientDisplacement _\SequenceIndex \longrightarrow0.
\]
This proves continuity of
$\ModelError $
in \eqref{eq:global-error}.

\ProofStep{The interpolating homotopy.}\label{step:global-homotopy}
\paraid{p-35ac6ff9ee97}
We interpolate the original metric and the model pseudometric to construct a continuous map
$\ApproximationHomotopy\colon\GHSpace\times\UnitInterval\to\GHSpace$
with the endpoint identities in \eqref{eq:global-domination-1} and \eqref{eq:global-domination-2}
and the estimate in \eqref{eq:global-track-control}.

\ProofSubstep{Definition and compactness.}\label{step:global-homotopy-definition}
\paraid{p-132d78a0f66f}
For
$\MetricSpace{\ComparisonCarrier}{\ComparisonMetric}\in\GHSpace$
and
$\HomotopyTime\in\UnitInterval$,
put
\begin{equation}\label{eq:interpolating-pseudometrics-1}
\InterpolatingPseudometric_{\ComparisonCarrier,\HomotopyTime}
 =(1-\HomotopyTime)\ComparisonMetric
       +\HomotopyTime\ModelPseudometric_\ComparisonCarrier,
\end{equation}
and
\begin{equation}\label{eq:interpolating-pseudometrics-2}
\ApproximationHomotopy(\MetricSpace{\ComparisonCarrier}{\ComparisonMetric},\HomotopyTime)
 =\MetricQuotient{\ComparisonCarrier}
       {\InterpolatingPseudometric_{\ComparisonCarrier,\HomotopyTime}}.
\end{equation}
Both
$\ComparisonMetric$
and
$\ModelPseudometric_\ComparisonCarrier$
are continuous pseudometrics on the compact space
$\ComparisonCarrier$.
By the definition,
$\InterpolatingPseudometric_{\ComparisonCarrier,\HomotopyTime}$
is a metric
generating the same topology on
$Y$
when
$\HomotopyTime<1$.

\paraid{p-7cd5915365f0}
By \eqref{eq:global-pseudometric-sum}
and the invariance under \eqref{eq:frame-action},
$\ModelPseudometric_\ComparisonCarrier$
is independent of the coordinate pairs chosen in the local models.
Moreover,
by \eqref{eq:input-isometry-1} and \eqref{eq:input-isometry-2},
every isometry
$\InputIsometry\colon\MetricSpace{\ComparisonCarrier}{\ComparisonMetric}
 \to\MetricSpace{\ComparisonCarrier_0}{\ComparisonMetric_0}$
satisfies
\[
 \ModelPseudometric_{\ComparisonCarrier_0}
   (\InputIsometry(\ComparisonPoint),\InputIsometry(\IntegrationPoint))
 =\ModelPseudometric_\ComparisonCarrier(\ComparisonPoint,\IntegrationPoint)
 \qquad(\ComparisonPoint,\IntegrationPoint\in\ComparisonCarrier).
\]
Thus
$\ApproximationHomotopy$
is well-defined on
$\GHSpace\times\UnitInterval$.

\ProofSubstep{The two endpoints.}\label{step:global-homotopy-endpoints}
\paraid{p-76f372559d9c}
We next prove \eqref{eq:global-domination-1} and \eqref{eq:global-domination-2}.
By \eqref{eq:interpolating-pseudometrics-1},
we obtain
$\InterpolatingPseudometric_{\ComparisonCarrier,0}=\ComparisonMetric$,
and
$\InterpolatingPseudometric_{\ComparisonCarrier,1}
       =\ModelPseudometric_\ComparisonCarrier$.
At the second endpoint,
for every
$\ComparisonPoint,\IntegrationPoint\in\ComparisonCarrier$,
\eqref{eq:global-pseudometric}
implies
that
$\ModelPseudometric_\ComparisonCarrier(\ComparisonPoint,\IntegrationPoint)=0$
if and only if
$\JointFeatureMap_\ComparisonCarrier(\ComparisonPoint)
       =\JointFeatureMap_\ComparisonCarrier(\IntegrationPoint)$.
Since
$\CompactFeatureImage_\ComparisonCarrier
 =\JointFeatureMap_\ComparisonCarrier(\ComparisonCarrier)$,
the induced map
\begin{equation}\label{eq:quotient-feature-isometry}
 \MetricQuotient{\ComparisonCarrier}{\ModelPseudometric_\ComparisonCarrier}
 \to\MetricSpace{\CompactFeatureImage_\ComparisonCarrier}
       {\RestrictedMetric{\BanachMetric}{\CompactFeatureImage_\ComparisonCarrier}}
\end{equation}
is defined as follows.
For each
$\ComparisonPoint\in\ComparisonCarrier$,
send
$[\ComparisonPoint]$
to
$\JointFeatureMap_\ComparisonCarrier(\ComparisonPoint)$.
It is a bijection.
For every
$\ComparisonPoint,\IntegrationPoint\in\ComparisonCarrier$,
the identity
\[
 \ModelPseudometric_\ComparisonCarrier(\ComparisonPoint,\IntegrationPoint)
 =\norm{\JointFeatureMap_\ComparisonCarrier(\ComparisonPoint)
        -\JointFeatureMap_\ComparisonCarrier(\IntegrationPoint)}_\BanachAmbientSpace
\]
shows that this map is an isometry.
By \eqref{eq:interpolating-pseudometrics-2},
the isometry
\eqref{eq:quotient-feature-isometry},
and the definitions of
$\ApproximationMap$
and
$\RealizationMap$,
we obtain the following equalities in
$\GHSpace$.
\[
 \ApproximationHomotopy(\MetricSpace{\ComparisonCarrier}{\ComparisonMetric},0)
  =\MetricSpace{\ComparisonCarrier}{\ComparisonMetric},
\]
and
\[
 \begin{aligned}
 \ApproximationHomotopy(\MetricSpace{\ComparisonCarrier}{\ComparisonMetric},1)
 &=\MetricQuotient{\ComparisonCarrier}{\ModelPseudometric_\ComparisonCarrier}=\MetricSpace{\CompactFeatureImage_\ComparisonCarrier}
       {\RestrictedMetric{\BanachMetric}{\CompactFeatureImage_\ComparisonCarrier}}
  =\RealizationMap\ApproximationMap\MetricSpace{\ComparisonCarrier}{\ComparisonMetric}.
 \end{aligned}
\]

\ProofSubstep{Variation of the time parameter.}\label{step:global-homotopy-time}
\paraid{p-8853a418623b}
Fix
$\MetricSpace{\ComparisonCarrier}{\ComparisonMetric}$.
To prove \eqref{eq:global-track-control},
let
$\FirstTime,\HomotopyTime\in\UnitInterval$.
By \eqref{eq:global-error},
\[
 \InterpolatingPseudometric_{\ComparisonCarrier,\FirstTime}
       -\InterpolatingPseudometric_{\ComparisonCarrier,\HomotopyTime}
 =(\FirstTime-\HomotopyTime)
       (\ModelPseudometric_\ComparisonCarrier-\ComparisonMetric),
\]
and
\[
 \norm{\InterpolatingPseudometric_{\ComparisonCarrier,\FirstTime}
       -\InterpolatingPseudometric_{\ComparisonCarrier,\HomotopyTime}}_\infty
 =|\FirstTime-\HomotopyTime|
       \ModelError\MetricSpace{\ComparisonCarrier}{\ComparisonMetric}.
\]
Applying \eqref{eq:pseudometric-bound} on the common carrier
$\ComparisonCarrier$,
we obtain
\[
 \begin{aligned}
 &\GHDistance\bigl(
   \ApproximationHomotopy(\MetricSpace{\ComparisonCarrier}{\ComparisonMetric},\FirstTime),
   \ApproximationHomotopy(\MetricSpace{\ComparisonCarrier}{\ComparisonMetric},\HomotopyTime)
   \bigr)\\
 &\qquad\leq\tfrac12
       \norm{\InterpolatingPseudometric_{\ComparisonCarrier,\FirstTime}
       -\InterpolatingPseudometric_{\ComparisonCarrier,\HomotopyTime}}_\infty
  =\tfrac12|\FirstTime-\HomotopyTime|
       \ModelError\MetricSpace{\ComparisonCarrier}{\ComparisonMetric}.
 \end{aligned}
\]
This is \eqref{eq:global-track-control}.
Figure~\ref{fig:pseudometric-homotopy} illustrates the interpolation
and the metric quotient at its endpoint.

\begin{figure}[H]
\centering
\begin{tikzpicture}[>=Stealth,font=\small]
\node at (0,2.15) {$\HomotopyTime=0$};
\node at (3.8,2.15) {$\HomotopyTime=\tfrac12$};
\node at (7.6,2.15) {$\HomotopyTime=1$};
\coordinate (leftzero) at (-0.375,0);
\coordinate (leftone) at (0.375,0);
\coordinate (lefttwo) at (0,1.4524);
\draw (leftzero) -- node[left] {$2$} (lefttwo)
 -- node[right] {$2$} (leftone)
 -- node[below] {$1$} (leftzero);
\foreach \point in {leftzero,leftone,lefttwo}
 \fill (\point) circle (1.4pt);
\node[below left] at (leftzero) {$\ComparisonPoint_0$};
\node[below right] at (leftone) {$\ComparisonPoint_1$};
\node[above] at (lefttwo) {$\ComparisonPoint_2$};
\begin{scope}[xshift=3.8cm]
 \coordinate (middlezero) at (-0.1875,0);
 \coordinate (middleone) at (0.1875,0);
 \coordinate (middletwo) at (0,1.4882);
 \draw (middlezero) -- node[left] {$2$} (middletwo)
 -- node[right] {$2$} (middleone)
 -- node[below] {$\tfrac12$} (middlezero);
 \foreach \point in {middlezero,middleone,middletwo}
  \fill (\point) circle (1.4pt);
 \node[below left] at (middlezero) {$\ComparisonPoint_0$};
 \node[below right] at (middleone) {$\ComparisonPoint_1$};
 \node[above] at (middletwo) {$\ComparisonPoint_2$};
\end{scope}
\begin{scope}[xshift=7.6cm]
 \draw (0,0) -- node[right] {$2$} (0,1.5);
 \fill (0,0) circle (1.4pt);
 \fill (0,1.5) circle (1.4pt);
 \node[below] at (0,0) {$[\ComparisonPoint_0]=[\ComparisonPoint_1]$};
 \node[above] at (0,1.5) {$[\ComparisonPoint_2]$};
\end{scope}
\draw[->] (1.15,0.7) -- (2.6,0.7);
\draw[->] (4.9,0.7) -- (6.4,0.7);
\node at (3.8,-1)
 {$\InterpolatingPseudometric_{\ComparisonCarrier,\HomotopyTime}
   =(1-\HomotopyTime)\ComparisonMetric
     +\HomotopyTime\ModelPseudometric_\ComparisonCarrier$};
\node at (3.8,-1.85)
 {$\begin{aligned}
  &\GHDistance\bigl(
   \ApproximationHomotopy(\MetricSpace{\ComparisonCarrier}{\ComparisonMetric},\FirstTime),
   \ApproximationHomotopy(\MetricSpace{\ComparisonCarrier}{\ComparisonMetric},\HomotopyTime)
   \bigr)\\
  &\qquad\leq\tfrac12|\FirstTime-\HomotopyTime|
       \ModelError\MetricSpace{\ComparisonCarrier}{\ComparisonMetric}
  \end{aligned}$};
\end{tikzpicture}
\caption{Interpolation on a three-point carrier.
The edge labels are distances for
$\InterpolatingPseudometric_{\ComparisonCarrier,\HomotopyTime}$.
The distance between
$\ComparisonPoint_0$
and
$\ComparisonPoint_1$
decreases from
$1$
to
$0$,
while their distances to
$\ComparisonPoint_2$
remain
$2$.
Distinct points can be identified at the endpoint.
The general track estimate \eqref{eq:global-track-control}
also holds there.}
\label{fig:pseudometric-homotopy}
\end{figure}


\ProofSubstep{Continuity when the space and time both vary.}\label{step:global-homotopy-continuity}
\paraid{p-f7c37d1cb847}
Let
\[
 \MetricSpace{\ComparisonCarrier_\SequenceIndex}{\ComparisonMetric_\SequenceIndex}
       \longrightarrow\MetricSpace{\ComparisonCarrier}{\ComparisonMetric}
       \quad\text{in }\GHSpace,
\]
and
\[
 \HomotopyTime_\SequenceIndex\longrightarrow\HomotopyTime
       \quad\text{in }\UnitInterval.
\]
We prove that
\[
 \GHDistance\bigl(
   \ApproximationHomotopy(
       \MetricSpace{\ComparisonCarrier_\SequenceIndex}{\ComparisonMetric_\SequenceIndex},
       \HomotopyTime_\SequenceIndex),
   \ApproximationHomotopy(\MetricSpace{\ComparisonCarrier}{\ComparisonMetric},\HomotopyTime)
   \bigr)\longrightarrow0.
\]
Use the correspondences
$\Correspondence_\SequenceIndex$
and the numbers
$\AmbientDisplacement_\SequenceIndex,\FeatureDisplacement_\SequenceIndex\to0$
from Step~\ref{step:global-continuity}.
First
we consider the case of
$\HomotopyTime_\SequenceIndex=\HomotopyTime$.
For
$(\ComparisonPoint_\SequenceIndex,\ComparisonPoint),
  (\IntegrationPoint_\SequenceIndex,\IntegrationPoint)\in\Correspondence_\SequenceIndex$,
the estimates in Step~\ref{step:global-continuity} imply
\[
 \begin{aligned}
 &\bigl|\InterpolatingPseudometric_{\ComparisonCarrier_\SequenceIndex,\HomotopyTime}
       (\ComparisonPoint_\SequenceIndex,\IntegrationPoint_\SequenceIndex)
       -\InterpolatingPseudometric_{\ComparisonCarrier,\HomotopyTime}
       (\ComparisonPoint,\IntegrationPoint)\bigr|\\
 &\qquad\leq(1-\HomotopyTime)
       |\ComparisonMetric_\SequenceIndex(\ComparisonPoint_\SequenceIndex,\IntegrationPoint_\SequenceIndex)
         -\ComparisonMetric(\ComparisonPoint,\IntegrationPoint)|
       +\HomotopyTime
       |\ModelPseudometric_{\ComparisonCarrier_\SequenceIndex}
          (\ComparisonPoint_\SequenceIndex,\IntegrationPoint_\SequenceIndex)
         -\ModelPseudometric_\ComparisonCarrier(\ComparisonPoint,\IntegrationPoint)|\\
 &\qquad\leq2(1-\HomotopyTime)\AmbientDisplacement_\SequenceIndex
       +2\HomotopyTime\FeatureDisplacement_\SequenceIndex.
 \end{aligned}
\]
At
$\HomotopyTime=1$,
the induced pseudometrics may identify distinct points,
and we use their metric quotients.
\Cref{lem:pseudometric-comparison} therefore applies also at this endpoint and yields
\begin{equation}\label{eq:fixed-time-homotopy-convergence}
 \begin{aligned}
 &\GHDistance\bigl(
   \ApproximationHomotopy(
       \MetricSpace{\ComparisonCarrier_\SequenceIndex}{\ComparisonMetric_\SequenceIndex},\HomotopyTime),
   \ApproximationHomotopy(\MetricSpace{\ComparisonCarrier}{\ComparisonMetric},\HomotopyTime)
   \bigr)\leq(1-\HomotopyTime)\AmbientDisplacement_\SequenceIndex
       +\HomotopyTime\FeatureDisplacement_\SequenceIndex\longrightarrow0.
 \end{aligned}
\end{equation}

\paraid{p-6a7506bb68e5}
We combine the fixed-time convergence with the track estimate to compare the values at
$\HomotopyTime_\SequenceIndex$
and
$\HomotopyTime$.
The continuity
of
$\ModelError$
proved in Step~\ref{step:global-continuity} implies
\[
 \ModelError\MetricSpace{\ComparisonCarrier_\SequenceIndex}{\ComparisonMetric_\SequenceIndex}
       \longrightarrow\ModelError\MetricSpace{\ComparisonCarrier}{\ComparisonMetric}.
\]
By the triangle inequality and \eqref{eq:global-track-control},
\[
 \begin{aligned}
 &\GHDistance\bigl(
   \ApproximationHomotopy(
       \MetricSpace{\ComparisonCarrier_\SequenceIndex}{\ComparisonMetric_\SequenceIndex},
       \HomotopyTime_\SequenceIndex),
   \ApproximationHomotopy(\MetricSpace{\ComparisonCarrier}{\ComparisonMetric},\HomotopyTime)
   \bigr)\\
 &\qquad\leq\GHDistance\bigl(
   \ApproximationHomotopy(
       \MetricSpace{\ComparisonCarrier_\SequenceIndex}{\ComparisonMetric_\SequenceIndex},
       \HomotopyTime_\SequenceIndex),
   \ApproximationHomotopy(
       \MetricSpace{\ComparisonCarrier_\SequenceIndex}{\ComparisonMetric_\SequenceIndex},\HomotopyTime)
   \bigr)+\GHDistance\bigl(
   \ApproximationHomotopy(
       \MetricSpace{\ComparisonCarrier_\SequenceIndex}{\ComparisonMetric_\SequenceIndex},\HomotopyTime),
   \ApproximationHomotopy(\MetricSpace{\ComparisonCarrier}{\ComparisonMetric},\HomotopyTime)
   \bigr)\\
 &\qquad\leq\tfrac12|\HomotopyTime_\SequenceIndex-\HomotopyTime|
       \ModelError\MetricSpace{\ComparisonCarrier_\SequenceIndex}{\ComparisonMetric_\SequenceIndex}
       +(1-\HomotopyTime)\AmbientDisplacement_\SequenceIndex
       +\HomotopyTime\FeatureDisplacement_\SequenceIndex\longrightarrow0.
 \end{aligned}
\]
Since
$\GHSpace\times\UnitInterval$
is metrizable,
this proves continuity of
$\ApproximationHomotopy\colon\GHSpace\times\UnitInterval\to\GHSpace$.
\end{proof}
